\chapter{Lay and Web Text}
\label{chap:bt-web}

\noindent
\begin{center}
\begin{tikzpicture}
\node[draw=ThesisNeutral, fill=ThesisPaper, rounded corners=3pt, line width=0.5pt,
      text width=0.84\textwidth, inner sep=12pt, align=justify,
      font=\footnotesize, text=ThesisInk] {
\textit{L'infarctus du myocarde (IDM) est la mort (nécrose) d'une zone plus ou
moins étendue du muscle cardiaque (myocarde). Les cellules musculaires cardiaques
de ce territoire ne parviennent plus à se contracter par manque d'apport en
oxygène et meurent en quelques heures.}
};
\end{tikzpicture}
\captionof{figure}{A general-public health page about myocardial infarction.}
\label{fig:bt-web}
\end{center}
\medskip

% === P1 — HOOK ===
\Cref{fig:bt-web} is the same heart attack once more, now explained to the person
who might have one. Every technical word is unpacked in passing, \emph{nécrose} as
the death of tissue, \emph{myocarde} as the heart muscle, and the reader is assumed
to know nothing. Text like this is the most abundant medical writing of all: there
is far more health information written for patients than there are journal
articles, let alone clinical notes. It is plainly not a clinical record. But it is
worth a look, because it shows what happens to medical language when it reaches the
public, and in particular what happens to names.

% === P2 — Fleck ===
That a fact changes as it travels is an old observation. In 1935 the physician
Ludwik Fleck argued that knowledge is produced inside a community of specialists
and is transformed as it moves outward to a wider audience
\citep{fleck_genesis_1979}. The cautious \emph{it appears that} of the laboratory
becomes the public's \emph{doctors have shown}; the careful technical term becomes
an everyday word. Writing for patients is not watered-down science, then, but a
different layer of the same knowledge, with its own form.

% === P3 — naming varies (Wüster foil, Cabré, Gaudin) ===
What changes most, for our purpose, is the naming. The founding ideal of
terminology, set out by the engineer Eugen Wüster, was one concept and one term,
each thing with a single correct name \citep{wuster_einfuhrung_1979}. Real usage
does not behave this way. The linguist María Teresa Cabré showed that a single
concept is named differently depending on who is speaking and why, and that this
variation is normal rather than a defect \citep{cabre_theories_2003}; in French,
Gaudin's socioterminology makes the same point about how terms spread from
specialists to the public \citep{gaudin_socioterminologie_2003}. The heart attack
of the clinical note, \emph{infarctus du myocarde}, coded \emph{I21}, is the same
event the patient calls a \emph{crise cardiaque}. One concept, several names, and
the public register invents the most of them.

% === P4 — UMLS, measured ===
Medicine has not resolved this multiplicity so much as catalogued it. The Unified
Medical Language System, described by Bodenreider, gathers more than sixty
vocabularies into a single map that links names to concepts
\citep{bodenreider_unified_2004}. The map is many-to-many: in its 2004 release it
held about 2.5 million names for some 900{,}000 concepts. The field's central
resource does not impose one name per concept, it records every name there is. The
dream of one concept and one term did not come true, and the infrastructure we
built around medical language is the admission that it did not.

% === P5 — Hacking + register, brief ===
Naming in medicine is not even a neutral act. As Ian Hacking observed, classifying
people can change the people classified, so a label can reshape the very thing it
names \citep{hacking_social_1999}. And the registers sit measurably apart: the
language of patients leans on pronouns and everyday words, where clinical and
scientific writing pile up nouns and exact terms \citep{biber_variation_1988}.

% === P6 — synthesis + bridge (closes the part) ===
The three texts complete a picture. The same medicine is written three ways, and
they differ not only in how much there is, from the rare clinical note to the
abundant article to the still more abundant page for patients, but in form, in
reasoning, and above all in how they name things. The text we can gather most
easily sits at the opposite end of this range from the text we actually need.
Using public data for the clinic is therefore not a question of quantity alone; it
means crossing from one register to another, and resolving the many names of a
thing into the one the clinic records. That is the problem the rest of this thesis
takes up.
